\documentclass[../main.tex]{subfiles}

\begin{document}

\justifying

\subsection{Descripción general de los experimentos}

Se busca caracterizar el proceso de formación de la red.
Hay 4 parámetros principales para variar, factor de empaquetamiento \(\phi\), temperatura \(T\), \% de Cross-Linkers \(\%\mathrm{CL}\) y longitud de potencial \(r_{c}\).
Primero probaremos los siguientes factores de empaquetamiento \(\phi\in\{0.01,0.03,0.05\}\) a una temperatura \(T=0.05\), \% de Cross-linkers \(\%\mathrm{CL}=5\%\) y longitud de potencial constanstante \(r_{c}=1.5 D_p\).
Después de explorar los factores de empaquetamiento se escojerá un factor de empaquetamiento y se explorará un rango de temperaturas \(T\in\{0.025, 0.035, 0.045, 0.055, 0.065, 0.075, 0.085, 0.095\}\).
Finalmente se explorará porcentajes de Cross-Linkers, \(\%\mathrm{CL}\in\{1\%,5\%,10\%\}\).

Lo que se va a analizar de la red es lo siguiente,
factor de estructura, 
tamaño de poro,
longtud de cadena promedio.
Se busca generar las siguientes gráficas,
tamaño del cluster más grande con respecto al tiempo,
la distribución del factor de estructura con respecto al tiempo,
distribución del tamaño de poro.





\end{document}
